\clearpage
\onecolumn
\section{Reproducibility, Artifact, and Media Manifest}
\label{sec:supp-reproducibility}

The supplement is built from tracked paper sources and project evidence
artifacts. The review-packet view below summarizes the visual artifact classes
used by the PDF: paired render atlases, material diagnostics, target-view
grounding renders, and selected navigation media.

\begin{center}
\centering
\includegraphics[width=0.98\textwidth,height=0.78\textheight,keepaspectratio]{figures/fig_supplement_review_packet_contact_sheet.png}
\captionsetup{hypcap=false}
\captionof{figure}{Review-packet media manifest for the supplement. The contact
sheet is a visual index of tracked render figures and selected still panels.}
\label{fig:supp-review-packet-contact-sheet}
\end{center}

\clearpage
\twocolumn

The main build target remains separate from the supplement build target. The
supplement PDF can be rebuilt with:

\begin{verbatim}
make -C paper acl27-supplement
\end{verbatim}

The current scaffold uses shared tables, shared figures, and ACL-local
supplement sections. Video material is represented by selected still panels and
manifest references. Raw frame directories, raw source scenes, scratch noMDL
USD trees, local model checkpoints, and private runtime logs remain outside the
safe default review packet.

\subsection{Primary Supplement Sources}

The readable supplement uses the following source classes:

\begin{itemize}
\item metric tables under \texttt{paper/shared/tables/};
\item render-atlas and scene-context panels under \texttt{paper/shared/figures/};
\item material-effect panels under \texttt{paper/shared/figures/};
\item InternNav selected still panels under \texttt{paper/shared/figures/};
\item generated per-case crop panels under \texttt{paper/shared/figures/supplement/};
\item media and video manifests under the project evidence tree.
\end{itemize}

The media manifests map selected cases to route status, failure reason, and
metrics. These manifests are used for traceability. They are not a license to
upload raw videos automatically.

\subsection{Anonymization Policy}

The supplement must satisfy the same anonymous-review scope as the main
paper. It must not include author names, local absolute paths, usernames,
private repository links, acknowledgments, or non-anonymous external services.
Any future media ZIP must be scanned separately because binary files can carry
filenames, metadata, or visible UI overlays that are not present in the PDF.

\subsection{Final Upload Decision}

The safe default is to upload the main PDF and OpenReview copy sources only.
The supplement PDF and media ZIP can be added after a final route and media
approval decision. If the supplement is uploaded, the author should also update
the private author checklist to record that optional media and supplement
contents were reviewed, anonymized, and approved.

\begin{figure*}[p]
\centering
\includegraphics[width=\textwidth,height=0.72\textheight,keepaspectratio]{figures/fig_supplement_source_boundary_companion.png}
\caption{Source-scope review companion. Render,
material, grounding, and navigation panels show the readable source classes
while the right panel marks excluded raw/private artifacts.}
\label{fig:supp-source-scope-companion}
\end{figure*}
\clearpage
